\chapter{通信：为另一个过程生成什么}
\label{ch:communication}
\chapterguide{消息的作用由接收者仍需区分的情境决定。生成、公开与采用分别发生，三者共同决定证据怎样进入组织。}

\section{局部历史中的通信选择}
Agent A已经检查了一部分证据，Agent B正决定是否继续某项方案。A可以转发已有观察，也可以专门生成一份与B的备选解释有关的似然摘要；后者可能在准备前尚未存在。

令$h_t^A$为A实际知道的历史。通信规则$\Gamma_A(h_t^A)$选择接收者、消息生成操作，以及立即公开、继续准备或等待。生成新片段仍是一条$\II_e\xrightarrow{F_a}v$；使片段对B可见是另外的事件。B的模式决定是否采用它，收到猜想不等于把猜想当作有效前提。

控制者的全局实验表不能代替A的局部信息。A可以依据公开计划或B发来的请求了解其需求。无论规则由谁生成，它在运行时都只能使用已经实际可见的材料。

\section{通信怎样改变诊断实验}
设某次实验的原片段为$v$，消息按与真实情境无关的核$Q(m\mid v)$生成。接收者看见的实验为
\[
 \widetilde K_a(m\mid\theta)=\sum_vQ(m\mid v)K_a(v\mid\theta).
\]
这是统计实验的后处理\citep{blackwell1953}。消息记录必须包含接收者能观察到的全部选择信号，适用时也包括发送标记与时刻。

\begin{theorem}[通信压缩的覆盖单调性]
\label{thm:garbling}
对每个实验和情境对，$h(\widetilde K_a^\theta,\widetilde K_a^{\theta'})\le h(K_a^\theta,K_a^{\theta'})$。因此，在相同实验分配集合下，$\widetilde\kap\le\kap$；不可辨别类只能合并或保持。
\end{theorem}
\begin{proof}
对每个消息$m$，Cauchy--Schwarz给出
\[
 \sqrt{\sum_vQ(m\mid v)P(v)\sum_vQ(m\mid v)P'(v)}
 \ge\sum_vQ(m\mid v)\sqrt{P(v)P'(v)}.
\]
对$m$求和，亲和度不会减小，区别量不会增大。逐项取加权和、最小值与最大值得到覆盖结论；原来相同的分布经过同样后处理仍相同。
\end{proof}

这不意味着压缩无用：传递与解码成本可能显著下降。定理提供必须一起核对的另一面——节约是否毁掉了任务依赖的区别。

\section{什么压缩完整保留了证据}
先考虑严格正概率表和确定性编码$M=\phi(V)$。某个编码单元是产生同一消息的片段集合。

\begin{proposition}[亲和度无损压缩]
\label{prop:sufficiency}
对一对情境，压缩前后亲和度相等，当且仅当每个编码单元中$K_a(v\mid\theta)/K_a(v\mid\theta')$为常数。对全部情境对无损，等价于每个单元中的似然向量$\bigl(K_a(v\mid\vartheta)\bigr)_{\vartheta\in\Theta}$互相成正比例。
\end{proposition}
\begin{proof}
确定性编码下，每个消息对应一个求和单元。上一定理的每个Cauchy--Schwarz差非负，和为零恰当各单元取等。严格正条件下，取等恰当平方根概率向量成比例，等价于所写似然比条件。对全部情境对应用该条件得到向量表述。
\end{proof}

无损压缩因此可以从接收者的备选模型出发设计。表达形式可以变化，只要这些似然区别被保留，统计证据就没有减少。保留某个任务当前决策所需的信息，有时还允许更强压缩；需要另行指定损失、先验与后续问题。

\section{为接收者生成一份充分记录}
允许实验安排根据已有观察自适应改变，但选择规则只能使用记录和情境独立随机性。$J$个成员在汇总点提交一份完整前缀的计数：每项携带情境信息的观察都恰好计入一个成员，不能遗漏用于决定下一实验的私有观察。第$j$个成员统计
\[
 C_{j,a,v}=\#\{t:\text{该调用属于成员 }j,\ A_t=a,V_t=v\}.
\]
汇总者计算$\ln L_n(\vartheta)=\sum_{j,a,v}C_{j,a,v}\ln K_a(v\mid\vartheta)$，便与接收全部原始记录得到相同的情境似然比和最大似然输出。

\begin{proposition}[计数消息的精确组织实现]
\label{prop:counts}
上述完整计数保留任意合法自适应组织的全部情境似然比。若总调用数不超过$n$，其一个固定编码使用至多
\[
 J|\calA||\VV|\lceil\log_2(n+1)\rceil
\]
bit，另加约定的传输封装成本。
\end{proposition}
\begin{proof}
任意完整记录的概率可写成$q_\Pi(T)\prod_{j,a,v}K_a(v\mid\vartheta)^{C_{j,a,v}}$，其中$q_\Pi(T)$是策略在该路径上的选择概率乘积，与$\vartheta$无关。将同一计数$C$的记录求和，得到
\[
 \PP_\vartheta(C)=\left(\sum_{T:\,C(T)=C}q_\Pi(T)\right)
       \prod_{j,a,v}K_a(v\mid\vartheta)^{C_{j,a,v}}.
\]
括号因子对全部情境相同，在似然比中消去。每个计数为$0$到$n$的整数，固定编码各用$\lceil\log_2(n+1)\rceil$ bit足够。
\end{proof}

准备这样的片段使接收者不必重放整个观察块。自适应选择本身不会破坏完整计数的充分性。困难发生在只统计被筛选公开的观察、漏掉决定筛选的私有证据时；此时应把缺少的证据计入，或推导包含选择与沉默的消息似然。充分性也不自动保证任意调度具有良好覆盖率，可靠性上界仍需另外满足相应分配条件。

\section{何时证据已经足够}
考虑接收者拥有完整证据前缀的情形。概率核严格正，选实验的规则可以自适应。令$\ell_t(\vartheta)=\ln L_t(\vartheta)$，阈值$b=\ln(m/\delta)$。首次存在候选$\vartheta$满足
\begin{equation}
 \ell_t(\vartheta)-\ell_t(\vartheta')\ge b
 \quad\text{对全部 }\vartheta'\in\Alt(\vartheta)
 \label{eq:stop}
\end{equation}
时输出$g(\vartheta)$并停止。规则可以在固定批次消息到达后才检查。

\begin{theorem}[随时有效的停止规则]
\label{thm:stop}
式\eqref{eq:stop}在任意实验选择策略下，作出错误停止的概率不超过$\delta$。若持续使用覆盖率$u>0$的独立分配，并在每个前缀检查，则停止时刻$\tau$几乎必然有限，且
\[
 \PP_\theta(\tau>n)\le\min\{1,me^{b/2}(1-u)^n\}.
\]
\end{theorem}
\begin{proof}
固定真实$\theta$。对每个冲突情境$\theta'$，$L_t(\theta')/L_t(\theta)$是初值1的非负鞅。若错误选中$\theta'$，该比值必曾达到$e^b$。非负鞅最大不等式给出概率至多$e^{-b}$；对至多$m$个冲突情境求和为$\delta$。

对独立分配，若在$n$时所有真实情境对冲突情境的对数似然差都达到$b$，必已停止。任意冲突情境违反该条件的概率，由平方根似然比的Markov不等式，至多$e^{b/2}(1-u)^n$。求和得到尾界，故停止几乎必然发生。
\end{proof}

随时有效的界通过鞅而非反复使用固定样本界成立\citep{howard2020}。仅仅在很多时刻重复检查某个固定置信阈值，通常不能自动保持总错误预算。

\section{沉默与选择性公开}
若A知道一个公平bit，双方约定可靠时隙，A仅在bit为1时发送标记，则B从“标记或沉默”完整恢复一bit。平均只有半次显式发送，但可见记录仍有一bit熵。发送次数、消息长度和完整记录信息量必须分别计算。

更一般地，以$S$表示发送与否，$V$在未发送时取固定空值，则
\[
 I_2(Y;S,V\mid Z)=I_2(Y;S\mid Z)
  +\PP(S=1)I_2(Y;V\mid Z,S=1).
\]
$Z$是接收者已有材料。等式由条件链式法则得出。忽略第一项，就会遗漏选择性公开本身的证据。若B不知道A是否已经完成推理，沉默的含义也随之改变。

\section{消息信息与任务损失}
接收者已有$Z$，需要预测目标$Y$，消息为$M$。最小平均对数损失由$H_2(Y\mid Z)$降为$H_2(Y\mid M,Z)$，差值为$I_2(Y;M\mid Z)$；条件交叉熵分解直接证明这一点。

若目标字母表大小为$q\ge2$，重构错误率至多$\delta\le1-1/q$，则条件Fano界给出
\[
 I_2(Y;M\mid Z)\ge H_2(Y\mid Z)-h_2(\delta)-\delta\log_2(q-1).
\]
完整消息用$L$个固定二元槽位表达时，左边至多$L$。信息与编码的基本联系源于通信理论\citep{shannon1948}；这里的诊断覆盖率进一步分析消息是否保留所有目标冲突对的区别。

\section{局部请求怎样改变内容}
若B只需知道某个候选是否被当前证据排除，A可以生成相应似然差；若B还要在许多候选之间选择，则完整似然向量或计数可能更适合。接收者的请求改变待保留的区别，从而改变应该执行的消息生成操作。

公开时机可以以式\eqref{eq:stop}所需的证据为依据，也可以按固定预算分批。更一般的通信价格、延迟与私有状态联合最优控制超出这个停止定理；但任何方案都必须把可见的沉默、筛选和时间信号放进同一个观察模型。

\begin{exercise}为什么将实验输出全部压缩成一个固定符号，会使任何诊断覆盖率变为零？\end{exercise}
\begin{exercise}给出一个选择性发送协议，说明只保留已发送内容的似然而忽略发送规则，会算错接收者后验。\end{exercise}
